\newlabel{supp:data}{{S1}{1}{Data}{section.1}{}}
\newlabel{supp:data-table}{{S1.1}{1}{Construction of the cell-week table}{subsection.1.1}{}}
\newlabel{supp:data-cov}{{S1.2}{2}{Covariates and missing values}{subsection.1.2}{}}
\newlabel{tab:covariates}{{S1}{3}{Covariates, resolution and source}{table.1}{}}
\newlabel{supp:methods-more}{{S2}{3}{Further methods}{section.2}{}}
\newlabel{supp:methods-fv}{{S2.1}{3}{Finite-volume discretisation and sea barriers}{subsection.2.1}{}}
\newlabel{eq:Qs}{{S2.1}{4}{Finite-volume discretisation and sea barriers}{equation.2.1}{}}
\newlabel{rem:boundary}{{1}{4}{}{remark.1}{}}
\newlabel{supp:methods-st}{{S2.2}{5}{Lattices and boundaries of the space-time field}{subsection.2.2}{}}
\newlabel{supp:methods-joint}{{S2.3}{5}{Joint model: identification and related models}{subsection.2.3}{}}
\newlabel{supp:methods-comp}{{S2.4}{6}{Computation}{subsection.2.4}{}}
\newlabel{sec:methods-persistence}{{S2.5}{7}{State dependence versus persistent heterogeneity}{subsection.2.5}{}}
\newlabel{supp:sim-more}{{S3}{8}{Further simulation studies}{section.3}{}}
\newlabel{sec:sim-dgp}{{S3.1}{8}{Data-generating process}{subsection.3.1}{}}
\newlabel{eq:sim}{{S3.1}{8}{Data-generating process}{equation.3.1}{}}
\newlabel{tab:s0}{{S2}{9}{S0: accuracy of the finite-volume SPDE field and of the empirical Bayes Laplace approximation, with fields drawn from the continuous Mat\'ern process (50 replicates per setting). Mean over replicates, with the SD across replicates in parentheses, except for the coverage of $b_1$, where it is the Monte Carlo standard error}{table.2}{}}
\newlabel{sec:sim-engine}{{S3.2}{9}{S0: accuracy of the field representation and of the Laplace engine}{subsection.3.2}{}}
\newlabel{sec:sim-persistence}{{S3.3}{9}{S1: state dependence versus persistent heterogeneity}{subsection.3.3}{}}
\newlabel{supp:app-more}{{S4}{10}{Further application results}{section.4}{}}
\newlabel{tab:s3}{{S3}{11}{S1: estimation of the state-dependence parameter $\delta $ (20 replicates per scenario)}{table.3}{}}
\newlabel{sec:app-data}{{S4.1}{11}{Reconstructing the cell-week table}{subsection.4.1}{}}
\newlabel{supp:app-jointpred}{{S4.2}{11}{Joint prediction when both outcomes are observed}{subsection.4.2}{}}
\newlabel{fig:gain-map}{{S1}{13}{Gain in joint log score per cell-week of the joint model over the product of the univariate static models, by country, on the 2023--2024 test period (blue: the joint model predicts better). Boxes: the five $10^\circ \times 10^\circ $ windows with the largest total gain on the 2022 validation year, selected without the test period; solid: the first}{figure.1}{}}
\newlabel{tab:blocks}{{S4}{13}{Latent blocks of the latent Gaussian model of Section~3.1 of the main paper, their precision matrices and hyperparameters}{table.4}{}}
\newlabel{supp:tables}{{S5}{13}{Additional tables}{section.5}{}}
\newlabel{tab:barrier}{{S5}{14}{Boundary treatments on a $60 \times 30$ lattice of $0.5^\circ $ cells with a sea strait of 2, 4 or 6 cells (about 110, 220, 330\,km): marginal variance of a coastal and of an inland land cell (nominal 1) and correlation between the two coastal cells facing each other across the strait, from 4{,}000 exact draws with $\rho = 400$\,km. ``Land only'' removes the strait from the lattice; ``sea = land'' keeps it as ordinary cells without data. The strait lies inside a closed rectangle. This is not the geometry of the application, which has an open coast, a buffer of sea cells and a reflecting outer edge, and the fitted ranges there are both above and below 400\,km, so the table shows the behaviour of the barrier only qualitatively}{table.5}{}}
\newlabel{tab:s5}{{S6}{14}{S2: full results for the escalation effect $\kUC $ of recent unrest on conflict, by scenario and estimator, over 20 replicates per scenario. Mean $\hkUC $, its empirical SD across replicates and the mean of the estimated SEs; coverage of the nominal 95\% intervals for $\kUC $ with its Monte Carlo SE; proportion of replicates whose interval excludes zero; mean estimate of the own-lag coefficient $d_c$ (truth 0.3); mean and SD of the estimated loading $\lambda $ of the shared space-time field (truth 0.8, joint model only); and mean log score, AUC-PR and onset AUC-PR of the conflict forecasts over the 26 test weeks. Onset rows are those of cells without a conflict event in the previous 52 weeks. The last two scenarios add unrest-only episodes (AR(1) with coefficient 0.95, Mat\'ern innovations of range 200\,km, marginal SD 0.6), which the joint model with one shared field must absorb into that field}{table.6}{}}
\newlabel{tab:app-joint-full}{{S7}{15}{Joint model of conflict and unrest, by region (Section~5 of the main paper): all hyperparameters and lag coefficients (posterior mean, with SE in parentheses) of the two univariate conflict models and of the joint model, fitted on 2018--2022. Subscripts C and U denote conflict and unrest. In the univariate model with $u + e + w + v$, the space-time field $v$ belongs to conflict alone; in the joint model it is shared, with loading $\lambda $ for unrest. Histories are $\log (1 + \text {events in the previous four weeks})$; neighbour histories are means over the rook neighbours, standardised on the training rows. The log marginal likelihood of the joint model includes the unrest data and is not comparable with those of the univariate models. Test scores are for 2023--2024 in the same region. The field $v$ is defined on the observed $1^\circ $ cells only, with reflecting edges, which departs from the boundary rule of Remark~1 of the main paper. SEs are plug-in values; in S1 and S2 such SEs are too small when a space-time field is fitted (ratio of empirical SD to mean SE up to 1.66 in S1 and 1.20 in S2 without escalation), so intervals from these models are likely too narrow. With a single shared space-time field, unrest-only episodes are absorbed into $v$, which can bias $\kUC $ towards zero; in S2 they inflated $\hat \lambda $ instead. The extended model has an additional unrest-specific space-time field $v_{\mathrm {U}}$. Test scores here use the probit approximation to $E\{\operatorname {expit}(\eta )\}$; the model comparisons in the text use common Monte Carlo draws}{table.7}{}}
\newlabel{tab:s3-full}{{S8}{16}{S1: estimation of the state-dependence parameter $\delta $ for all five truths and all step lengths of the space-time field $v$ (20 replicates per truth). The two truths with slow $v$ have an AR(1) coefficient of 0.97 per week instead of 0.9. SD/SE is the ratio of the empirical SD of $\hat \delta $ to the mean estimated SE; values above 1 mean that the plug-in intervals are too narrow. Time is the mean fit time per replicate}{table.8}{}}
